\documentclass[a4paper]{article}
\usepackage{amsfonts}
\usepackage{amsmath}
\usepackage{amssymb}
\usepackage{graphicx}     

\begin{document}
  
\noindent \large Auteur: Dara van Engelen\\
\noindent \large Studierichting: Informatica\\
\noindent \large Oefeningenreeks 2E nr. 26.7\\

\medskip

\normalsize

$(x+1)(x+7)\log _6 2 + (x+2)(x+6)\log _6 3 = x - \log _3 2 \cdot \log _6 243$\\

$\Leftrightarrow (x+1)(x+7)\log _6 2 + (x+2)(x+6)\log _6 3 = x - \dfrac{\log _6 2}{\log _6 3} \cdot \log _6 243$\\

$\Leftrightarrow (x+1)(x+7)\log _6 2 + (x+2)(x+6)\log _6 3 = x - \dfrac{\log _6 2}{\log _6 3} \cdot \log _6 3^5$\\

$\Leftrightarrow (x+1)(x+7)\log _6 2 + (x+2)(x+6)\log _6 3 = x - \dfrac{\log _6 2}{\log _6 3} \cdot 5\log _6 3$\\

$\Leftrightarrow (x^2+8x+7)\log _6 2 + (x^2+8x+12)\log _6 3 = x - 5 \log _6 2$\\

$\Leftrightarrow (x^2+8x+7)\log _6 2 + (x^2+8x+12)\log _6 3  - x +  5 \log _6 2 = 0 $\\

$\Leftrightarrow (x^2+8x+12)\log _6 2 + (x^2+8x+12)\log _6 3  - x = 0 $\\

$\Leftrightarrow (x^2 + 8x +12)(\log _6 2 + \log _6 3) - x = 0$\\

$\Leftrightarrow (x^2 + 8x +12)\log _6 6 - x = 0$\\

$\Leftrightarrow x^2+7x+12 = 0 $\\

$D = 49 - 4\cdot1\cdot12 = 1$\\

$x_1 = \dfrac{-7+1}{2} = -3$\\

$x_2 = \dfrac{-7-1}{2} = -4$\\

\begin{thebibliography}{999}
%\bibitem{a} Referentie
\end{thebibliography}
\end{document} 